\documentclass[../DoAn.tex]{subfiles}
\begin{document}

Trong kỷ nguyên số hóa hiện nay, việc quản lý công việc và cộng tác nhóm hiệu quả đã trở thành yếu tố then chốt quyết định sự thành bại của các dự án công nghệ cũng như các hoạt động học thuật. Chương 1 này đóng vai trò giới thiệu tổng quan về đề tài xây dựng ứng dụng quản lý công việc và cộng tác nhóm HustFlow. Chương này nêu bật sự cần thiết của đề tài, đặt nền móng lý do phát triển hệ thống và giới thiệu cấu trúc chung của báo cáo đồ án tốt nghiệp. Nội dung chi tiết của chương được chia thành bốn đề mục chính bao gồm: (i) đặt vấn đề thực tế trong quản lý công việc nhóm ở phần 1.1, (ii) xác định mục tiêu và phạm vi nghiên cứu của đề tài ở phần 1.2, (iii) đề xuất định hướng giải pháp kỹ thuật ở phần 1.3, và (iv) trình bày bố cục chi tiết của các chương còn lại trong đồ án ở phần 1.4.

\section{Đặt vấn đề}
\label{section:1.1}

Trong bối cảnh phát triển mạnh mẽ của công nghệ thông tin và sự thay đổi trong mô hình làm việc hiện đại, các tổ chức và nhóm phát triển ngày càng có nhu cầu cao về việc tối ưu hóa quy trình làm việc nhóm. Việc phối hợp nhịp nhàng giữa các thành viên, phân chia nhiệm vụ rõ ràng và theo dõi tiến độ một cách trực quan là các yếu tố quyết định để đảm bảo dự án hoàn thành đúng thời hạn và đạt chất lượng yêu cầu. Việc thiếu đi các công cụ hỗ trợ quản lý hiệu quả thường dẫn đến sự lãng phí tài nguyên, chồng chéo công việc và suy giảm hiệu suất làm việc chung của toàn tổ chức.

Khảo sát thực tế tại các nhóm làm việc, đặc biệt là các nhóm phát triển phần mềm và các nhóm nghiên cứu học thuật, cho thấy nhiều khó khăn trong quy trình phối hợp hàng ngày. Các vấn đề cốt lõi bao gồm: (i) việc theo dõi trạng thái công việc gặp nhiều trở ngại khi chỉ sử dụng danh sách tác vụ tĩnh, dẫn đến khó khăn trong việc nhận diện các điểm nghẽn; (ii) quản lý mốc thời gian và thời hạn hoàn thành phức tạp khi các công việc con và thời hạn chung bị tách rời; (iii) sự thiếu kiểm soát đối với mối quan hệ phụ thuộc giữa các thẻ công việc dẫn đến tình trạng hoàn thành sai thứ tự quy trình; (iv) hiện tượng không đồng bộ dữ liệu thời gian thực giữa các thành viên cùng thao tác, dễ gây ra xung đột thông tin; và (v) gánh nặng hành chính khi người quản trị phải tổng hợp số liệu tiến độ và phân bổ tài nguyên một cách thủ công.

Từ những thách thức nêu trên, việc phát triển một hệ thống quản lý công việc tích hợp các chế độ xem đa chiều, đồng bộ thời gian thực và hỗ trợ tự động hóa bằng trí tuệ nhân tạo là một yêu cầu cấp thiết. Giải quyết tốt bài toán này sẽ giúp nâng cao năng lực cộng tác của các thành viên, giảm thiểu thời gian quản lý thủ công và hỗ trợ ra quyết định nhanh chóng cho người quản lý. Hệ thống này không chỉ phục vụ đắc lực cho các nhóm phát triển công nghệ mà còn có thể áp dụng rộng rãi trong nhiều lĩnh vực hoạt động nhóm khác nhau.

\section{Mục tiêu và phạm vi đề tài}
\label{section:1.2}

Hiện nay, trên thị trường đã có một số công cụ quản lý công việc phổ biến như Trello, Jira hay Asana. Mỗi hệ thống đều có những ưu điểm riêng biệt nhằm phục vụ các đối tượng người dùng cụ thể. Ứng dụng Trello nổi bật với giao diện bảng Kanban đơn giản, dễ tiếp cận nhưng thiếu các góc nhìn tiến độ chuyên sâu như biểu đồ Gantt hay lịch biểu tích hợp sẵn. Ngược lại, hệ thống Jira cung cấp các quy trình quản lý vô cùng chặt chẽ cho các dự án phần mềm lớn nhưng lại có giao diện phức tạp, đòi hỏi chi phí bản quyền cao và khó tiếp cận đối với các nhóm quy mô nhỏ hoặc các nhóm sinh viên. Trong khi đó, ứng dụng Asana hỗ trợ giao diện dòng thời gian tương đối trực quan nhưng chưa tối ưu khả năng cập nhật thời gian thực khi nhiều thành viên đồng thời thao tác trên cùng một màn hình.

Bên cạnh đó, phần lớn các công cụ hiện tại chưa tích hợp sâu khả năng của trí tuệ nhân tạo thế hệ mới để hỗ trợ người dùng trong quá trình soạn thảo và kiểm soát chất lượng tác vụ. Người sử dụng vẫn phải thực hiện thủ công các công việc lặp đi lặp lại như viết mô tả chi tiết, đề xuất danh mục công việc con hay tổng hợp báo cáo sức khỏe dự án. Khoảng trống này tạo ra cơ hội để xây dựng một giải pháp quản lý công việc thế hệ mới kết hợp hài hòa giữa tính trực quan của Kanban, tính bao quát của biểu đồ dòng thời gian, khả năng tương tác thời gian thực và sự trợ giúp thông minh từ AI.

Mục tiêu cụ thể của đề tài này là thiết kế và xây dựng ứng dụng HustFlow nhằm giải quyết các hạn chế nêu trên. Các mục tiêu kỹ thuật và nghiệp vụ chính bao gồm: (i) xây dựng bảng Kanban kéo thả linh hoạt tích hợp phân quyền vai trò chi tiết; (ii) phát triển các góc nhìn tiến độ đa chiều gồm Lịch biểu, biểu đồ dòng thời gian Gantt và chế độ xem chia đôi màn hình; (iii) thiết lập cơ chế kiểm soát và cảnh báo mối quan hệ phụ thuộc giữa các công việc để tránh vòng lặp; (iv) tích hợp trợ lý trí tuệ nhân tạo để tự động hóa việc khởi tạo mô tả, gợi ý danh sách công việc con và tổng hợp báo cáo tiến độ; và (v) cài đặt hệ thống giới hạn tài nguyên và thanh toán để quản lý việc nâng cấp gói cước.

Phạm vi nghiên cứu và triển khai của đề tài tập trung vào việc phát triển một ứng dụng web hoàn chỉnh, có khả năng đáp ứng nhu cầu cộng tác của các nhóm phát triển phần mềm, các nhóm dự án học thuật và các nhà quản lý dự án quy mô vừa và nhỏ. Hệ thống hỗ trợ mô hình đa tổ chức làm việc, trong đó mỗi tổ chức có không gian làm việc riêng biệt, được bảo mật và đồng bộ hóa thời gian thực qua kết nối WebSocket. Các chi tiết thiết kế giao diện, cấu trúc cơ sở dữ liệu và kịch bản thực nghiệm hệ thống sẽ được phân tích rõ hơn ở các chương sau.

\section{Định hướng giải pháp}
\label{section:1.3}

Để hiện thực hóa các mục tiêu đã đề ra, đề tài định hướng xây dựng hệ thống HustFlow dựa trên kiến trúc ứng dụng web hiện đại. Công nghệ cốt lõi được lựa chọn bao gồm: (i) Next.js App Router kết hợp thư viện React nhằm tối ưu hiệu năng hiển thị và mang lại trải nghiệm người dùng mượt mà; (ii) hệ quản trị cơ sở dữ liệu MySQL cùng Prisma ORM đảm bảo tính toàn vẹn và an toàn dữ liệu; (iii) dịch vụ Pusher làm nền tảng truyền tải dữ liệu thời gian thực qua giao thức WebSocket; (iv) Clerk quản lý quá trình xác thực người dùng và phân cấp tổ chức làm việc; và (v) OpenAI API cung cấp các khả năng xử lý ngôn ngữ tự nhiên cho trợ lý AI. Sự kết hợp này nhằm xây dựng một hệ thống ổn định, có độ trễ thấp và khả năng mở rộng tốt.

Giải pháp kỹ thuật của HustFlow tổ chức thông tin công việc theo cấu trúc phân cấp chặt chẽ từ không gian tổ chức đến các bảng công việc, các cột danh sách và các thẻ tác vụ chi tiết. Ở khía cạnh quản lý trạng thái ứng dụng, hệ thống sử dụng React Query để quản lý và lưu bộ nhớ đệm dữ liệu phía máy khách, kết hợp với các sự kiện Pusher được phát đi mỗi khi có thay đổi dữ liệu trên máy chủ. Khi một thành viên thực hiện thao tác kéo thả thẻ, chỉnh sửa thuộc tính hoặc gửi bình luận, máy chủ sẽ xử lý thay đổi và gửi tín hiệu đồng bộ. Các máy khách khác đang tham gia cùng bảng sẽ nhận được tín hiệu này để tự động cập nhật bộ nhớ đệm mà không cần tải lại toàn bộ trang web.

Đồ án tốt nghiệp đã đạt được các kết quả thực tiễn quan trọng bao gồm: (i) xây dựng thành công ứng dụng web HustFlow đáp ứng đầy đủ các yêu cầu quản lý công việc đa chiều; (ii) phát triển thuật toán phát hiện vòng lặp phụ thuộc hiệu quả và cơ chế cảnh báo thông minh; (iii) xây dựng giải pháp trợ lý AI tương tác trực tiếp giúp rút ngắn thời gian chuẩn bị dữ liệu và phân tích báo cáo; và (iv) tích hợp hoàn chỉnh hệ thống thanh toán Stripe để quản lý vòng đời gói cước của tổ chức. Hệ thống được đánh giá đạt các tiêu chuẩn về tính dễ dùng, độ tin cậy và hiệu năng phản hồi nhanh chóng trong môi trường hoạt động thực tế.

\section{Bố cục đồ án}
\label{section:1.4}

Báo cáo đồ án tốt nghiệp này ngoài chương mở đầu được cấu trúc thành năm chương tiếp theo như dưới đây.

Chương 2 tập trung vào phần Khảo sát hiện trạng và Phân tích yêu cầu hệ thống. Chương này đánh giá chi tiết các quy trình quản lý công việc thực tế và phân tích các ưu nhược điểm của các giải pháp hiện có, từ đó xây dựng các yêu cầu chức năng và phi chức năng cho hệ thống HustFlow. Các yêu cầu chức năng được mô hình hóa chi tiết thông qua biểu đồ ca sử dụng tổng quát và các biểu đồ phân rã ca sử dụng, kèm theo đặc tả luồng sự kiện cho các chức năng cốt lõi của ứng dụng.

Chương 3 trình bày về Nền tảng lý thuyết và Công nghệ sử dụng. Chương này giới thiệu chi tiết cơ sở khoa học của các công nghệ được áp dụng như mô hình Next.js App Router, thư viện React, hệ quản trị cơ sở dữ liệu MySQL, Prisma ORM, cơ chế Pusher và API của OpenAI. Đồng thời, nội dung chương cũng đưa ra các lập luận so sánh kỹ thuật để minh chứng cho sự phù hợp của các công nghệ đã chọn đối với yêu cầu bài toán.

Chương 4 mô tả chi tiết quá trình Phân tích thiết kế, Triển khai và Đánh giá hệ thống. Nội dung chương bao gồm thiết kế kiến trúc hệ thống tổng thể, thiết kế biểu đồ gói, biểu đồ tuần tự cho các luồng nghiệp vụ chính, sơ đồ quan hệ thực thể cơ sở dữ liệu và thống kê mã nguồn dự án. Chương này cũng trình bày các kịch bản kiểm thử chức năng, kiểm thử phi chức năng cùng các kết quả đo đạc hiệu năng thực tế.

Chương 5 giới thiệu chi tiết Các giải pháp và Đóng góp nổi bật của đồ án. Đây là chương cốt lõi làm rõ giá trị sáng tạo kỹ thuật của đề tài, bao gồm giải pháp đồng bộ dữ liệu thời gian thực tối ưu kết hợp Pusher và React Query, thuật toán kiểm tra tính phụ thuộc vòng lặp giữa các thẻ công việc, giải pháp tích hợp trợ lý trí tuệ nhân tạo để tự động hóa báo cáo và sinh nội dung, và giải pháp phân quyền vai trò truy cập kết hợp Stripe billing.

Chương 6 đưa ra Kết luận và Hướng phát triển của đề tài. Chương này tổng kết lại những kết quả đạt được, tự đánh giá các ưu điểm và hạn chế còn tồn tại của hệ thống HustFlow. Từ đó, đề tài đề xuất định hướng nghiên cứu tiếp theo và các giải pháp cải tiến công nghệ nhằm nâng cao hiệu suất và trải nghiệm người dùng trong tương lai.

Như vậy, chương này đã phác thảo những nét khái quát nhất về đề tài nghiên cứu và xây dựng hệ thống HustFlow. Thông qua việc phân tích bối cảnh thực tế và các khó khăn hiện hữu trong hoạt động phối hợp nhóm, đề tài đã định vị rõ ràng mục tiêu xây dựng một ứng dụng trực quan, tích hợp sâu các chế độ xem đa chiều và hỗ trợ thông minh từ trí tuệ nhân tạo. Định hướng giải pháp kỹ thuật dựa trên các công nghệ hiện đại cũng đã được lựa chọn nhằm đảm bảo khả năng vận hành tối ưu. Những nội dung tổng quan này sẽ làm cơ sở để tiến hành phân tích chi tiết các yêu cầu nghiệp vụ và khảo sát hiện trạng thực tế, nội dung này sẽ được trình bày cụ thể trong chương tiếp theo – Chương 2.

\end{document}